\section{嘉道危机可跳转事件节点}
以下锚点将战争、河工、旗务、边贸、禁烟和条约放入同一可导航的时间框架。文书能证明的行政动作与需要地方材料核验的社会后果在每项后分别说明。
\subsection{1790—1799}
\eventanchor{EQ-1796-JIAQING-ACCESSION}{嘉庆帝正式即位，白莲教战争与财政…}嘉庆元年（1796年），嘉庆帝正式即位，白莲教战争与财政压力成为初政核心问题。核心取证：《清仁宗实录》嘉庆元年相关条；宫中朱批奏折、内阁大库题本与《清史稿》相关志传。此类文件确认机构作出或收到的行政动作，不能跳过地方执行与反馈。来源保留：以嘉庆元年（1796年）的同期文书为定位起点；官修记录、奏报或外交文本服务特定机构立场，具体日期、规模、地方执行与对手视角须以独立材料复核。
\eventanchor{EQ-1797-WHITE-LOTUS-WAR}{白莲教战争持续扩展，清军依赖乡勇…}嘉庆二年（1797年），白莲教战争持续扩展，清军依赖乡勇、团练与地方资源。核心取证：《清仁宗实录》嘉庆二年相关条；军机处档、战事方略及地方奏报。编年证词定位国家叙述中的时间，却需同地方或对方材料校验范围。来源保留：以嘉庆二年（1797年）的同期文书为定位起点；官修记录、奏报或外交文本服务特定机构立场，具体日期、规模、地方执行与对手视角须以独立材料复核。
\eventanchor{EQ-1798-MILITIA-MOBILIZATION}{朝廷加强对乡勇、团练和地方绅士资…}嘉庆三年（1798年），朝廷加强对乡勇、团练和地方绅士资源的动员以应对白莲教战争。核心取证：《清仁宗实录》嘉庆三年相关条；地方志、督抚奏折及地方档案。编年证词定位国家叙述中的时间，却需同地方或对方材料校验范围。来源保留：以嘉庆三年（1798年）的同期文书为定位起点；官修记录、奏报或外交文本服务特定机构立场，具体日期、规模、地方执行与对手视角须以独立材料复核。
\eventanchor{EQ-1799-HESHEN-PURGE}{乾隆帝去世后，嘉庆帝逮治和珅并清…}嘉庆四年（1799年），乾隆帝去世后，嘉庆帝逮治和珅并清理其政治网络。核心取证：《清仁宗实录》嘉庆四年相关条；宫中朱批奏折、内阁大库题本与《清史稿》相关志传。此类文件确认机构作出或收到的行政动作，不能跳过地方执行与反馈。来源保留：以嘉庆四年（1799年）的同期文书为定位起点；官修记录、奏报或外交文本服务特定机构立场，具体日期、规模、地方执行与对手视角须以独立材料复核。
\subsection{1800—1809}
\eventanchor{EQ-1800-WHITE-LOTUS-FINANCE}{白莲教战争军费、捐输与地方筹饷加…}嘉庆五年（1800年），白莲教战争军费、捐输与地方筹饷加重国家财政压力。核心取证：《清仁宗实录》嘉庆五年相关条；户部奏销、盐政、河工或海关档案。此类文件确认机构作出或收到的行政动作，不能跳过地方执行与反馈。来源保留：以嘉庆五年（1800年）的同期文书为定位起点；官修记录、奏报或外交文本服务特定机构立场，具体日期、规模、地方执行与对手视角须以独立材料复核。
\eventanchor{EQ-1801-WHITE-LOTUS-SUPPRESSION}{清军在若干战区逐步压缩白莲教起事…}嘉庆六年（1801年），清军在若干战区逐步压缩白莲教起事者活动空间。核心取证：《清仁宗实录》嘉庆六年相关条；军机处档、战事方略及地方奏报。编年证词定位国家叙述中的时间，却需同地方或对方材料校验范围。来源保留：以嘉庆六年（1801年）的同期文书为定位起点；官修记录、奏报或外交文本服务特定机构立场，具体日期、规模、地方执行与对手视角须以独立材料复核。
\eventanchor{EQ-1802-WAR-REFUGEE-RELIEF}{战区流民、复垦和赈济成为地方善后…}嘉庆七年（1802年），战区流民、复垦和赈济成为地方善后事务。核心取证：《清仁宗实录》嘉庆七年相关条；地方志、督抚奏折及地方档案。编年证词定位国家叙述中的时间，却需同地方或对方材料校验范围。来源保留：以嘉庆七年（1802年）的同期文书为定位起点；官修记录、奏报或外交文本服务特定机构立场，具体日期、规模、地方执行与对手视角须以独立材料复核。
\eventanchor{EQ-1803-FORBIDDEN-CITY-ATTACK}{陈德等闯入紫禁城，暴露京师警卫与…}嘉庆八年（1803年），陈德等闯入紫禁城，暴露京师警卫与内廷管理问题。核心取证：《清仁宗实录》嘉庆八年相关条；宫中朱批奏折、内阁大库题本与《清史稿》相关志传。此类文件确认机构作出或收到的行政动作，不能跳过地方执行与反馈。来源保留：以嘉庆八年（1803年）的同期文书为定位起点；官修记录、奏报或外交文本服务特定机构立场，具体日期、规模、地方执行与对手视角须以独立材料复核。
\eventanchor{EQ-1804-COASTAL-PIRACY}{东南海盗活动加剧，福建、广东沿海…}嘉庆九年（1804年），东南海盗活动加剧，福建、广东沿海军政压力上升。核心取证：《清仁宗实录》嘉庆九年相关条；军机处档、战事方略及地方奏报。编年证词定位国家叙述中的时间，却需同地方或对方材料校验范围。来源保留：以嘉庆九年（1804年）的同期文书为定位起点；官修记录、奏报或外交文本服务特定机构立场，具体日期、规模、地方执行与对手视角须以独立材料复核。
\eventanchor{EQ-1805-WHITE-LOTUS-WAR-END}{白莲教大规模战争基本结束，清廷转…}嘉庆十年（1805年），白莲教大规模战争基本结束，清廷转入善后和治安重整。核心取证：《清仁宗实录》嘉庆十年相关条；军机处档、战事方略及地方奏报。编年证词定位国家叙述中的时间，却需同地方或对方材料校验范围。来源保留：以嘉庆十年（1805年）的同期文书为定位起点；官修记录、奏报或外交文本服务特定机构立场，具体日期、规模、地方执行与对手视角须以独立材料复核。
\eventanchor{EQ-1806-TSAI-QIAN-PIRACY}{蔡牵等海盗集团持续活动，清廷调整…}嘉庆十一年（1806年），蔡牵等海盗集团持续活动，清廷调整水师与沿海防务。核心取证：《清仁宗实录》嘉庆十一年相关条；军机处档、战事方略及地方奏报。编年证词定位国家叙述中的时间，却需同地方或对方材料校验范围。来源保留：以嘉庆十一年（1806年）的同期文书为定位起点；官修记录、奏报或外交文本服务特定机构立场，具体日期、规模、地方执行与对手视角须以独立材料复核。
\eventanchor{EQ-1807-COASTAL-NAVAL-REFORM}{清廷继续通过水师调度、招抚和地方…}嘉庆十二年（1807年），清廷继续通过水师调度、招抚和地方协作应对海盗。核心取证：《清仁宗实录》嘉庆十二年相关条；宫中朱批奏折、内阁大库题本与《清史稿》相关志传。此类文件确认机构作出或收到的行政动作，不能跳过地方执行与反馈。来源保留：以嘉庆十二年（1807年）的同期文书为定位起点；官修记录、奏报或外交文本服务特定机构立场，具体日期、规模、地方执行与对手视角须以独立材料复核。
\eventanchor{EQ-1808-PIRATE-ALLIANCES}{郑一嫂、张保仔等海盗网络在华南海…}嘉庆十三年（1808年），郑一嫂、张保仔等海盗网络在华南海域形成强大联盟。核心取证：《清仁宗实录》嘉庆十三年相关条；地方志、督抚奏折及地方档案。编年证词定位国家叙述中的时间，却需同地方或对方材料校验范围。来源保留：以嘉庆十三年（1808年）的同期文书为定位起点；官修记录、奏报或外交文本服务特定机构立场，具体日期、规模、地方执行与对手视角须以独立材料复核。
\eventanchor{EQ-1809-SOUTH-CHINA-SEA-CAMPAIGN}{清军、地方势力与外来船只在华南海…}嘉庆十四年（1809年），清军、地方势力与外来船只在华南海域共同影响剿盗行动。核心取证：《清仁宗实录》嘉庆十四年相关条；军机处档、战事方略及地方奏报。编年证词定位国家叙述中的时间，却需同地方或对方材料校验范围。来源保留：以嘉庆十四年（1809年）的同期文书为定位起点；官修记录、奏报或外交文本服务特定机构立场，具体日期、规模、地方执行与对手视角须以独立材料复核。
\subsection{1810—1819}
\eventanchor{EQ-1810-ZHANG-BAOZAI-SURRENDER}{张保仔接受招安，华南大海盗联盟瓦…}嘉庆十五年（1810年），张保仔接受招安，华南大海盗联盟瓦解。核心取证：《清仁宗实录》嘉庆十五年相关条；军机处档、战事方略及地方奏报。编年证词定位国家叙述中的时间，却需同地方或对方材料校验范围。来源保留：以嘉庆十五年（1810年）的同期文书为定位起点；官修记录、奏报或外交文本服务特定机构立场，具体日期、规模、地方执行与对手视角须以独立材料复核。
\eventanchor{EQ-1811-BANNER-LIVELIHOOD}{八旗生计、债务和人口供养问题持续…}嘉庆十六年（1811年），八旗生计、债务和人口供养问题持续成为中央与地方的治理议题。核心取证：《清仁宗实录》嘉庆十六年相关条；地方志、督抚奏折及地方档案。编年证词定位国家叙述中的时间，却需同地方或对方材料校验范围。来源保留：以嘉庆十六年（1811年）的同期文书为定位起点；官修记录、奏报或外交文本服务特定机构立场，具体日期、规模、地方执行与对手视角须以独立材料复核。
\eventanchor{EQ-1812-TIANLI-PREPARATION}{天理教组织活动在直隶、河南、山东…}嘉庆十七年（1812年），天理教组织活动在直隶、河南、山东等地被官府逐步察觉。核心取证：《清仁宗实录》嘉庆十七年相关条；地方志、督抚奏折及地方档案。编年证词定位国家叙述中的时间，却需同地方或对方材料校验范围。来源保留：以嘉庆十七年（1812年）的同期文书为定位起点；官修记录、奏报或外交文本服务特定机构立场，具体日期、规模、地方执行与对手视角须以独立材料复核。
\eventanchor{EQ-1813-TIANLI-UPRISING}{天理教起事者攻入紫禁城部分区域，…}嘉庆十八年（1813年），天理教起事者攻入紫禁城部分区域，京师发生严重安全危机。核心取证：《清仁宗实录》嘉庆十八年相关条；军机处档、战事方略及地方奏报。编年证词定位国家叙述中的时间，却需同地方或对方材料校验范围。来源保留：以嘉庆十八年（1813年）的同期文书为定位起点；官修记录、奏报或外交文本服务特定机构立场，具体日期、规模、地方执行与对手视角须以独立材料复核。
\eventanchor{EQ-1814-TIANLI-SUPPRESSION}{天理教起事被镇压，相关人员遭逮捕…}嘉庆十九年（1814年），天理教起事被镇压，相关人员遭逮捕和处置。核心取证：《清仁宗实录》嘉庆十九年相关条；宫中朱批奏折、内阁大库题本与《清史稿》相关志传。此类文件确认机构作出或收到的行政动作，不能跳过地方执行与反馈。来源保留：以嘉庆十九年（1814年）的同期文书为定位起点；官修记录、奏报或外交文本服务特定机构立场，具体日期、规模、地方执行与对手视角须以独立材料复核。
\eventanchor{EQ-1815-GRAND-CANAL-REPAIRS}{漕运、河工和仓储问题在战后财政压…}嘉庆二十年（1815年），漕运、河工和仓储问题在战后财政压力下持续受到讨论。核心取证：《清仁宗实录》嘉庆二十年相关条；户部奏销、盐政、河工或海关档案。此类文件确认机构作出或收到的行政动作，不能跳过地方执行与反馈。来源保留：以嘉庆二十年（1815年）的同期文书为定位起点；官修记录、奏报或外交文本服务特定机构立场，具体日期、规模、地方执行与对手视角须以独立材料复核。
\eventanchor{EQ-1816-AMHERST-MISSION}{阿美士德使团来华，礼仪与外交交涉…}嘉庆二十一年（1816年），阿美士德使团来华，礼仪与外交交涉未获突破。核心取证：《清仁宗实录》嘉庆二十一年相关条；《筹办夷务始末》相关年卷与中外照会、条约文本。此类文件确认机构作出或收到的行政动作，不能跳过地方执行与反馈。来源保留：以嘉庆二十一年（1816年）的同期文书为定位起点；官修记录、奏报或外交文本服务特定机构立场，具体日期、规模、地方执行与对手视角须以独立材料复核。
\eventanchor{EQ-1817-SALT-ARREARS}{盐课、商欠和财政亏空继续困扰中央…}嘉庆二十二年（1817年），盐课、商欠和财政亏空继续困扰中央财政。核心取证：《清仁宗实录》嘉庆二十二年相关条；户部奏销、盐政、河工或海关档案。此类文件确认机构作出或收到的行政动作，不能跳过地方执行与反馈。来源保留：以嘉庆二十二年（1817年）的同期文书为定位起点；官修记录、奏报或外交文本服务特定机构立场，具体日期、规模、地方执行与对手视角须以独立材料复核。
\eventanchor{EQ-1818-FLOOD-RELIEF}{黄淮水患和赈务成为地方奏报与中央…}嘉庆二十三年（1818年），黄淮水患和赈务成为地方奏报与中央拨款的重要议题。核心取证：《清仁宗实录》嘉庆二十三年相关条；地方志、督抚奏折及地方档案。编年证词定位国家叙述中的时间，却需同地方或对方材料校验范围。来源保留：以嘉庆二十三年（1818年）的同期文书为定位起点；官修记录、奏报或外交文本服务特定机构立场，具体日期、规模、地方执行与对手视角须以独立材料复核。
\eventanchor{EQ-1819-OPIUM-TRADE-GROWTH}{鸦片贸易增长引起官员对白银流动、…}嘉庆二十四年（1819年），鸦片贸易增长引起官员对白银流动、海防和社会秩序的担忧。核心取证：《清仁宗实录》嘉庆二十四年相关条；户部奏销、盐政、河工或海关档案。此类文件确认机构作出或收到的行政动作，不能跳过地方执行与反馈。来源保留：以嘉庆二十四年（1819年）的同期文书为定位起点；官修记录、奏报或外交文本服务特定机构立场，具体日期、规模、地方执行与对手视角须以独立材料复核。
\subsection{1820—1829}
\eventanchor{EQ-1820-DAO-GUANG-ACCESSION}{道光帝即位，财政、河工、旗务和沿…}嘉庆二十五年（1820年），道光帝即位，财政、河工、旗务和沿海贸易问题进入新朝。核心取证：《清仁宗实录》嘉庆二十五年相关条；宫中朱批奏折、内阁大库题本与《清史稿》相关志传。此类文件确认机构作出或收到的行政动作，不能跳过地方执行与反馈。来源保留：以嘉庆二十五年（1820年）的同期文书为定位起点；官修记录、奏报或外交文本服务特定机构立场，具体日期、规模、地方执行与对手视角须以独立材料复核。
\eventanchor{EQ-1821-IMPERIAL-ECONOMY}{道光初年继续检讨内务府、官员经费…}道光元年（1821年），道光初年继续检讨内务府、官员经费和财政节用。核心取证：《清宣宗实录》道光元年相关条；户部奏销、盐政、河工或海关档案。此类文件确认机构作出或收到的行政动作，不能跳过地方执行与反馈。来源保留：以道光元年（1821年）的同期文书为定位起点；官修记录、奏报或外交文本服务特定机构立场，具体日期、规模、地方执行与对手视角须以独立材料复核。
\eventanchor{EQ-1822-YELLOW-RIVER-FLOOD}{黄河水患、堤工和灾赈问题反复进入…}道光二年（1822年），黄河水患、堤工和灾赈问题反复进入中央议程。核心取证：《清宣宗实录》道光二年相关条；地方志、督抚奏折及地方档案。编年证词定位国家叙述中的时间，却需同地方或对方材料校验范围。来源保留：以道光二年（1822年）的同期文书为定位起点；官修记录、奏报或外交文本服务特定机构立场，具体日期、规模、地方执行与对手视角须以独立材料复核。
\eventanchor{EQ-1823-BANNER-DEBT}{旗人债务、典当和生计问题持续引起…}道光三年（1823年），旗人债务、典当和生计问题持续引起官府讨论。核心取证：《清宣宗实录》道光三年相关条；地方志、督抚奏折及地方档案。编年证词定位国家叙述中的时间，却需同地方或对方材料校验范围。来源保留：以道光三年（1823年）的同期文书为定位起点；官修记录、奏报或外交文本服务特定机构立场，具体日期、规模、地方执行与对手视角须以独立材料复核。
\eventanchor{EQ-1824-GRAND-CANAL-CRISIS}{运河水运和漕粮转输面临河工、淤积…}道光四年（1824年），运河水运和漕粮转输面临河工、淤积与财政问题。核心取证：《清宣宗实录》道光四年相关条；户部奏销、盐政、河工或海关档案。此类文件确认机构作出或收到的行政动作，不能跳过地方执行与反馈。来源保留：以道光四年（1824年）的同期文书为定位起点；官修记录、奏报或外交文本服务特定机构立场，具体日期、规模、地方执行与对手视角须以独立材料复核。
\eventanchor{EQ-1825-SOUTHWEST-MINING-UNREST}{云贵矿业、移民和治安纠纷继续考验…}道光五年（1825年），云贵矿业、移民和治安纠纷继续考验地方治理。核心取证：《清宣宗实录》道光五年相关条；地方志、督抚奏折及地方档案。编年证词定位国家叙述中的时间，却需同地方或对方材料校验范围。来源保留：以道光五年（1825年）的同期文书为定位起点；官修记录、奏报或外交文本服务特定机构立场，具体日期、规模、地方执行与对手视角须以独立材料复核。
\eventanchor{EQ-1826-JAHANGIR-KHOJA-UPRISING}{张格尔和卓进入喀什噶尔，天山南路…}道光六年（1826年），张格尔和卓进入喀什噶尔，天山南路爆发大规模反清动荡。核心取证：《清宣宗实录》道光六年相关条；军机处档、战事方略及地方奏报。编年证词定位国家叙述中的时间，却需同地方或对方材料校验范围。来源保留：以道光六年（1826年）的同期文书为定位起点；官修记录、奏报或外交文本服务特定机构立场，具体日期、规模、地方执行与对手视角须以独立材料复核。
\eventanchor{EQ-1827-KASHGAR-RECOVERY}{清军反攻并收复喀什噶尔等地，张格…}道光七年（1827年），清军反攻并收复喀什噶尔等地，张格尔势力受挫。核心取证：《清宣宗实录》道光七年相关条；军机处档、战事方略及地方奏报。编年证词定位国家叙述中的时间，却需同地方或对方材料校验范围。来源保留：以道光七年（1827年）的同期文书为定位起点；官修记录、奏报或外交文本服务特定机构立场，具体日期、规模、地方执行与对手视角须以独立材料复核。
\eventanchor{EQ-1828-JAHANGIR-CAPTURE}{张格尔被捕并押解处置，清廷加强南…}道光八年（1828年），张格尔被捕并押解处置，清廷加强南疆军政控制。核心取证：《清宣宗实录》道光八年相关条；军机处档、理藩院档或相关方略。编年证词定位国家叙述中的时间，却需同地方或对方材料校验范围。来源保留：以道光八年（1828年）的同期文书为定位起点；官修记录、奏报或外交文本服务特定机构立场，具体日期、规模、地方执行与对手视角须以独立材料复核。
\eventanchor{EQ-1829-XINJIANG-FINANCE}{南疆战后军费、驻防与贸易秩序成为…}道光九年（1829年），南疆战后军费、驻防与贸易秩序成为清廷边务财政负担。核心取证：《清宣宗实录》道光九年相关条；户部奏销、盐政、河工或海关档案。此类文件确认机构作出或收到的行政动作，不能跳过地方执行与反馈。来源保留：以道光九年（1829年）的同期文书为定位起点；官修记录、奏报或外交文本服务特定机构立场，具体日期、规模、地方执行与对手视角须以独立材料复核。
\subsection{1830—1839}
\eventanchor{EQ-1830-RUSSO-QING-TRADE}{恰克图等边境贸易和俄清交涉持续影…}道光十年（1830年），恰克图等边境贸易和俄清交涉持续影响北部边务。核心取证：《清宣宗实录》道光十年相关条；《筹办夷务始末》相关年卷与中外照会、条约文本。此类文件确认机构作出或收到的行政动作，不能跳过地方执行与反馈。来源保留：以道光十年（1830年）的同期文书为定位起点；官修记录、奏报或外交文本服务特定机构立场，具体日期、规模、地方执行与对手视角须以独立材料复核。
\eventanchor{EQ-1831-CENTRAL-ASIA-CHOLERA}{疾病流行与交通网络对军队、城市和…}道光十一年（1831年），疾病流行与交通网络对军队、城市和商路造成压力。核心取证：《清宣宗实录》道光十一年相关条；地方志、督抚奏折及地方档案。编年证词定位国家叙述中的时间，却需同地方或对方材料校验范围。来源保留：以道光十一年（1831年）的同期文书为定位起点；官修记录、奏报或外交文本服务特定机构立场，具体日期、规模、地方执行与对手视角须以独立材料复核。
\eventanchor{EQ-1832-COASTAL-DEFENCE}{鸦片走私增长背景下的海防、缉私和…}道光十二年（1832年），鸦片走私增长背景下的海防、缉私和地方官责任成为争议。核心取证：《清宣宗实录》道光十二年相关条；宫中朱批奏折、内阁大库题本与《清史稿》相关志传。此类文件确认机构作出或收到的行政动作，不能跳过地方执行与反馈。来源保留：以道光十二年（1832年）的同期文书为定位起点；官修记录、奏报或外交文本服务特定机构立场，具体日期、规模、地方执行与对手视角须以独立材料复核。
\eventanchor{EQ-1833-OPIUM-POLICY-DEBATE}{官员围绕禁烟、征税、白银和海防提…}道光十三年（1833年），官员围绕禁烟、征税、白银和海防提出相互竞争的方案。核心取证：《清宣宗实录》道光十三年相关条；宫中朱批奏折、内阁大库题本与《清史稿》相关志传。此类文件确认机构作出或收到的行政动作，不能跳过地方执行与反馈。来源保留：以道光十三年（1833年）的同期文书为定位起点；官修记录、奏报或外交文本服务特定机构立场，具体日期、规模、地方执行与对手视角须以独立材料复核。
\eventanchor{EQ-1834-NAPIER-AFFAIR}{律劳卑抵广州后与两广官员发生外交…}道光十四年（1834年），律劳卑抵广州后与两广官员发生外交摩擦。核心取证：《清宣宗实录》道光十四年相关条；《筹办夷务始末》相关年卷与中外照会、条约文本。此类文件确认机构作出或收到的行政动作，不能跳过地方执行与反馈。来源保留：以道光十四年（1834年）的同期文书为定位起点；官修记录、奏报或外交文本服务特定机构立场，具体日期、规模、地方执行与对手视角须以独立材料复核。
\eventanchor{EQ-1835-GUANGZHOU-TRADE-REGULATION}{广州贸易、行商债务和外商管理继续…}道光十五年（1835年），广州贸易、行商债务和外商管理继续在地方与中央之间协商。核心取证：《清宣宗实录》道光十五年相关条；户部奏销、盐政、河工或海关档案。此类文件确认机构作出或收到的行政动作，不能跳过地方执行与反馈。来源保留：以道光十五年（1835年）的同期文书为定位起点；官修记录、奏报或外交文本服务特定机构立场，具体日期、规模、地方执行与对手视角须以独立材料复核。
\eventanchor{EQ-1836-OPIUM-LEGALIZATION-DEBATE}{许乃济等围绕鸦片弛禁、征税与禁绝…}道光十六年（1836年），许乃济等围绕鸦片弛禁、征税与禁绝提出政策论争。核心取证：《清宣宗实录》道光十六年相关条；宫中朱批奏折、内阁大库题本与《清史稿》相关志传。此类文件确认机构作出或收到的行政动作，不能跳过地方执行与反馈。来源保留：以道光十六年（1836年）的同期文书为定位起点；官修记录、奏报或外交文本服务特定机构立场，具体日期、规模、地方执行与对手视角须以独立材料复核。
\eventanchor{EQ-1837-CANTON-SMUGGLING}{广东缉私与鸦片走私之间的冲突持续…}道光十七年（1837年），广东缉私与鸦片走私之间的冲突持续升级。核心取证：《清宣宗实录》道光十七年相关条；地方志、督抚奏折及地方档案。编年证词定位国家叙述中的时间，却需同地方或对方材料校验范围。来源保留：以道光十七年（1837年）的同期文书为定位起点；官修记录、奏报或外交文本服务特定机构立场，具体日期、规模、地方执行与对手视角须以独立材料复核。
\eventanchor{EQ-1838-LIN-ZEXU-APPOINTMENT}{林则徐受命赴广东查禁鸦片，禁烟政…}道光十八年（1838年），林则徐受命赴广东查禁鸦片，禁烟政策转为高压执行。核心取证：《清宣宗实录》道光十八年相关条；宫中朱批奏折、内阁大库题本与《清史稿》相关志传。此类文件确认机构作出或收到的行政动作，不能跳过地方执行与反馈。来源保留：以道光十八年（1838年）的同期文书为定位起点；官修记录、奏报或外交文本服务特定机构立场，具体日期、规模、地方执行与对手视角须以独立材料复核。
\eventanchor{EQ-1839-HUMEN-OPIUM-DESTRUCTION}{林则徐在虎门销毁收缴鸦片，禁烟冲…}道光十九年（1839年），林则徐在虎门销毁收缴鸦片，禁烟冲突公开化。核心取证：《清宣宗实录》道光十九年相关条；《筹办夷务始末》相关年卷与中外照会、条约文本。此类文件确认机构作出或收到的行政动作，不能跳过地方执行与反馈。来源保留：以道光十九年（1839年）的同期文书为定位起点；官修记录、奏报或外交文本服务特定机构立场，具体日期、规模、地方执行与对手视角须以独立材料复核。
\subsection{1840—1849}
\eventanchor{EQ-1840-FIRST-OPIUM-WAR}{英国舰队来华，第一次鸦片战争爆发}道光二十年（1840年），英国舰队来华，第一次鸦片战争爆发。核心取证：《清宣宗实录》道光二十年相关条；军机处档、战事方略及地方奏报。编年证词定位国家叙述中的时间，却需同地方或对方材料校验范围。来源保留：以道光二十年（1840年）的同期文书为定位起点；官修记录、奏报或外交文本服务特定机构立场，具体日期、规模、地方执行与对手视角须以独立材料复核。
\eventanchor{EQ-1841-CHUANBI-NEGOTIATIONS}{穿鼻草约等谈判和沿海战事并行，清…}道光二十一年（1841年），穿鼻草约等谈判和沿海战事并行，清廷内部对和战意见分歧。核心取证：《清宣宗实录》道光二十一年相关条；《筹办夷务始末》相关年卷与中外照会、条约文本。此类文件确认机构作出或收到的行政动作，不能跳过地方执行与反馈。来源保留：以道光二十一年（1841年）的同期文书为定位起点；官修记录、奏报或外交文本服务特定机构立场，具体日期、规模、地方执行与对手视角须以独立材料复核。
\eventanchor{EQ-1841-CANTON-BATTLES}{广州周边战事和赔款谈判反复进行}道光二十一年（1841年），广州周边战事和赔款谈判反复进行。核心取证：《清宣宗实录》道光二十一年相关条；军机处档、战事方略及地方奏报。编年证词定位国家叙述中的时间，却需同地方或对方材料校验范围。来源保留：以道光二十一年（1841年）的同期文书为定位起点；官修记录、奏报或外交文本服务特定机构立场，具体日期、规模、地方执行与对手视角须以独立材料复核。
\eventanchor{EQ-1842-NANJING-TREATY}{《南京条约》签订，清廷割让香港并…}道光二十二年（1842年），《南京条约》签订，清廷割让香港并开放五口通商、支付赔款。核心取证：《清宣宗实录》道光二十二年相关条；《筹办夷务始末》相关年卷与中外照会、条约文本。此类文件确认机构作出或收到的行政动作，不能跳过地方执行与反馈。来源保留：以道光二十二年（1842年）的同期文书为定位起点；官修记录、奏报或外交文本服务特定机构立场，具体日期、规模、地方执行与对手视角须以独立材料复核。
\eventanchor{EQ-1843-BOGUE-TREATY}{《虎门条约》补充领事裁判、协定关…}道光二十三年（1843年），《虎门条约》补充领事裁判、协定关税等安排。核心取证：《清宣宗实录》道光二十三年相关条；《筹办夷务始末》相关年卷与中外照会、条约文本。此类文件确认机构作出或收到的行政动作，不能跳过地方执行与反馈。来源保留：以道光二十三年（1843年）的同期文书为定位起点；官修记录、奏报或外交文本服务特定机构立场，具体日期、规模、地方执行与对手视角须以独立材料复核。
\eventanchor{EQ-1844-WANGXIA-WHAMPOA-TREATIES}{《望厦条约》《黄埔条约》分别确立…}道光二十四年（1844年），《望厦条约》《黄埔条约》分别确立美法在华通商与相关权利。核心取证：《清宣宗实录》道光二十四年相关条；《筹办夷务始末》相关年卷与中外照会、条约文本。此类文件确认机构作出或收到的行政动作，不能跳过地方执行与反馈。来源保留：以道光二十四年（1844年）的同期文书为定位起点；官修记录、奏报或外交文本服务特定机构立场，具体日期、规模、地方执行与对手视角须以独立材料复核。
\eventanchor{EQ-1845-TREATY-PORT-OPENING}{五口通商制度进入具体开埠、税则和…}道光二十五年（1845年），五口通商制度进入具体开埠、税则和城市管理阶段。核心取证：《清宣宗实录》道光二十五年相关条；户部奏销、盐政、河工或海关档案。此类文件确认机构作出或收到的行政动作，不能跳过地方执行与反馈。来源保留：以道光二十五年（1845年）的同期文书为定位起点；官修记录、奏报或外交文本服务特定机构立场，具体日期、规模、地方执行与对手视角须以独立材料复核。
\eventanchor{EQ-1846-MISSIONARY-RETURN}{条约环境下传教活动与地方教案问题…}道光二十六年（1846年），条约环境下传教活动与地方教案问题逐渐增多。核心取证：《清宣宗实录》道光二十六年相关条；地方志、督抚奏折及地方档案。编年证词定位国家叙述中的时间，却需同地方或对方材料校验范围。来源保留：以道光二十六年（1846年）的同期文书为定位起点；官修记录、奏报或外交文本服务特定机构立场，具体日期、规模、地方执行与对手视角须以独立材料复核。
\eventanchor{EQ-1847-GUANGZHOU-CONFLICT}{广州周边的中英冲突显示条约执行和…}道光二十七年（1847年），广州周边的中英冲突显示条约执行和城市主权问题尚未解决。核心取证：《清宣宗实录》道光二十七年相关条；《筹办夷务始末》相关年卷与中外照会、条约文本。此类文件确认机构作出或收到的行政动作，不能跳过地方执行与反馈。来源保留：以道光二十七年（1847年）的同期文书为定位起点；官修记录、奏报或外交文本服务特定机构立场，具体日期、规模、地方执行与对手视角须以独立材料复核。
\eventanchor{EQ-1848-YANGTZE-TRADE}{通商口岸贸易网络向长江和沿海延伸…}道光二十八年（1848年），通商口岸贸易网络向长江和沿海延伸，地方商人开始适应新制度。核心取证：《清宣宗实录》道光二十八年相关条；户部奏销、盐政、河工或海关档案。此类文件确认机构作出或收到的行政动作，不能跳过地方执行与反馈。来源保留：以道光二十八年（1848年）的同期文书为定位起点；官修记录、奏报或外交文本服务特定机构立场，具体日期、规模、地方执行与对手视角须以独立材料复核。
\eventanchor{EQ-1849-GUANGZHOU-CITY-ENTRY}{英国要求进入广州城的争议再次激化…}道光二十九年（1849年），英国要求进入广州城的争议再次激化，地方民众与官府关系受影响。核心取证：《清宣宗实录》道光二十九年相关条；地方志、督抚奏折及地方档案。编年证词定位国家叙述中的时间，却需同地方或对方材料校验范围。来源保留：以道光二十九年（1849年）的同期文书为定位起点；官修记录、奏报或外交文本服务特定机构立场，具体日期、规模、地方执行与对手视角须以独立材料复核。
\subsection{1850—1859}
\eventanchor{EQ-1850-XIANFENG-ACCESSION}{咸丰帝即位，财政紧张、内乱和对外…}道光三十年（1850年），咸丰帝即位，财政紧张、内乱和对外压力同时加重。核心取证：《清宣宗实录》道光三十年相关条；宫中朱批奏折、内阁大库题本与《清史稿》相关志传。此类文件确认机构作出或收到的行政动作，不能跳过地方执行与反馈。来源保留：以道光三十年（1850年）的同期文书为定位起点；官修记录、奏报或外交文本服务特定机构立场，具体日期、规模、地方执行与对手视角须以独立材料复核。
